\subsection{Recurrence and scalar bounds for determinant potentials}
\label{app:determinant-details}

We prove the algebraic lemmas for the all-degree induction in
\cref{det:first-section}, using \(\mathcal C_m\), \(\mathcal R_r\), and
the binomial sums of
\eqref{det:cleared-polynomial}--\eqref{det:scalar-budget}.

\subsubsection*{Coefficient extraction for the recurrence}

\begin{proof}[Proof of \cref{det:stable-recurrence}]
For $q\geq1$, write
\[
 H_{st}^{(q)}=\sum_i y_i^q e_s(\widehat y_i)e_t(\widehat y_i).
\]
Here and below $e_j=0$ for $j<0$.  Direct substitution in
\eqref{det:cleared-polynomial} gives
\begin{align}
 \mathcal C_m=e_1(D_m^{\mathrm{sym}})^2\biggl[
 &2\sum_{s=1}^{m-1}s(s+1)\frac{e_{s+1}}{e_s}
 -2\sum_{s=1}^{m-1}\frac{H_{ss}^{(1)}}{e_s^2}
 -\sum_{s,t=1}^{m-1}\frac{H_{st}^{(1)}}{e_se_t}\notag\\
 &+\frac1{e_1}\sum_{s=1}^{m-1}\frac{H_{s0}^{(2)}}{e_s}
 -\frac{2e_2}{e_1}\biggr].
 \label{det:C-expanded}
\end{align}
This also shows that the denominators cancel.  The case $m=2$ gives
$\mathcal C_2=0$.

We calculate the coefficient of $e_{r+1}$ in
\eqref{det:C-expanded} for $m=r+1$.  Let
$\mathscr E(z)=\sum_{j=0}^{r+1}e_jz^j=\prod_i(1+y_i z)$.
Summing \(H_{st}^{(1)}z^sw^t\) gives the formal power series identity
\begin{align}
 \sum_{s,t\geq0}H_{st}^{(1)}z^sw^t
 &=\mathscr E(z)\mathscr E(w)
       \sum_i\frac{y_i}{(1+y_i z)(1+y_i w)}\notag\\
 &=\frac{z\mathscr E'(z)\mathscr E(w)
          -w\mathscr E'(w)\mathscr E(z)}{z-w}.
 \label{det:H-generating}
\end{align}
Extracting the coefficient of $z^sw^t$ in
\eqref{det:H-generating} gives, for $0\leq s,t\leq r$,
\[
 H_{st}^{(1)}=\sum_{j=0}^{\min(s,t)}
       (s+t+1-2j)e_j e_{s+t+1-j},
\]
where $e_j=0$ for $j>r+1$ in this calculation with $r+1$ variables.
Since $j\leq r$ in every summand, this expression is affine in
$e_{r+1}$. For $1\leq s,t\leq r$, differentiation with respect to
that formal variable therefore yields
\begin{equation}\label{det:H-coefficient}
 \frac{\partial H_{st}^{(1)}}{\partial e_{r+1}}
 =\begin{cases}
   (2r+1-s-t)e_{s+t-r},&s+t\geq r,\\
   0,&s+t<r.
  \end{cases}
\end{equation}
Similarly,
\[
 \sum_{s\geq0}H_{s0}^{(2)}z^s
 =\frac{e_1\mathscr E(z)-\mathscr E'(z)}z,
 \qquad
 H_{s0}^{(2)}=e_1e_{s+1}-(s+2)e_{s+2}.
\]
It follows that
\begin{equation}\label{det:H2-coefficient}
 \frac{\partial H_{s0}^{(2)}}{\partial e_{r+1}}
 =\begin{cases}
   0,&s\leq r-2,\\
   -(r+1),&s=r-1,\\
   e_1,&s=r.
  \end{cases}
\end{equation}
These formulas show that $\mathcal C_{r+1}$ depends affinely on
$e_{r+1}$.  Setting $e_{r+1}=0$ leaves $e_r^2\mathcal C_r$.
Indeed, append a zero variable to $r$ positive variables.  The terms
with $s=r$ vanish at every nonzero coordinate, and every term from the
zero coordinate has a factor $y_i$ or $y_i^2$.  This proves the claimed
restriction as a polynomial identity.

Since the ratio of $e_1(D_{r+1}^{\mathrm{sym}})^2$ to
$\mathcal P_r/e_{r+1}$ is $e_r$, equations
\eqref{det:H-coefficient}--\eqref{det:H2-coefficient} give
\begin{align}
 \frac{\mathcal R_r}{\mathcal P_r}
 ={}&2r(r+1)
 -2e_r\sum_{\substack{1\leq s\leq r\\2s\geq r}}
       \frac{(2r+1-2s)e_{2s-r}}{e_s^2}\notag\\
 &-e_r\sum_{\substack{1\leq s,t\leq r\\s+t\geq r}}
       \frac{(2r+1-s-t)e_{s+t-r}}{e_se_t}
 +1-(r+1)\frac{e_r}{e_1e_{r-1}}.
 \label{det:R-before-grouping}
\end{align}
The diagonal term $s=r$ contributes $-2$.  In the double sum, the
terms with $s=r$ or $t=r$ contribute $-(r^2+r-1)$.  Together with
$2r(r+1)+1$, these terms give the constant $r(r+1)$.

The remaining terms with $s+t=r$ have the form
$e_r/(e_pe_{r-p})$.  The double sum contributes the factor
$2-\mathbf1_{2p=r}$, and the diagonal sum contributes
$2\mathbf1_{2p=r}$.  The last term of
\eqref{det:R-before-grouping} adds $\mathbf1_{p=1}$.
Their total coefficient is the first coefficient in
\eqref{det:R-recurrence}.

For the terms with $s+t=r+k>r$ and $s,t<r$, choose $s\leq t$ and
write $s=k+d$, $t=r-d$.  Then
$1\leq d\leq\lfloor(r-1)/2\rfloor$ and $1\leq k\leq r-2d$.
The two orders in the double sum and the diagonal contribution give
$2+\mathbf1_{k=r-2d}$, while
$2r+1-s-t=r-k+1$.  This is the second sum in
\eqref{det:R-recurrence}, and proves the recurrence.
\end{proof}

\subsubsection*{The scalar binomial estimate}

\begin{proof}[Proof of \cref{det:binomial-budget}]
Let $S_{r,d}$ denote the inner sum in $J_r$, and put
\[
 t_{r,d}=\frac{\binom{r-d}{d}}{\binom rd},\qquad
 t_{r,d}=0\quad\text{if }2d>r.
\]
For $\ell=r-2d$, the identities
\[
 \sum_{k=1}^{\ell}\binom{k+d}{d}
     =\binom{r-d+1}{d+1}-1,\qquad
 \sum_{k=1}^{\ell}k\binom{k+d}{d}
     =(d+1)\binom{r-d+1}{d+2}
\]
give
\begin{align}
 S_{r,d}={}&\frac2{\binom rd}
  \left[(r+1)\left(\binom{r-d+1}{d+1}-1\right)
       -(d+1)\binom{r-d+1}{d+2}\right]\notag\\
 &+(2d+1)t_{r,d}.
 \label{det:Srd}
\end{align}
Also,
\[
 \frac{t_{r,d+1}}{t_{r,d}}
 =\frac{(r-2d)(r-2d-1)}{(r-d)^2}.
\]
We claim that
\begin{equation}\label{det:telescoping}
 S_{r,d}\leq\frac{2r^3}{r-1}
  \left(\frac{t_{r,d}}{d+1}-\frac{t_{r,d+1}}{d+2}\right).
\end{equation}
To check the sign explicitly, divide \eqref{det:Srd} by $t_{r,d}$.
After omitting the negative term
$-2(r+1)/\binom{r-d}{d}$, the result is
\[
 B_{r,d}=\frac{2(r-d+1)(r+2d^2+3d+2)}{(d+1)(d+2)}+2d+1.
\]
Put $u=d-1\geq0$ and $v=r-2d-1\geq0$.  Over the positive denominator
$(r-1)(d+1)(d+2)(r-d)^2$, the numerator of
\[
 \frac{2r^3}{r-1}\left(\frac1{d+1}
    -\frac{(r-2d)(r-2d-1)}{(d+2)(r-d)^2}\right)-B_{r,d}
\]
is
\begin{align*}
 {}&2uv^4+(4u^3+25u^2+29u)v^3\\
 &+(16u^4+102u^3+206u^2+140u+14)v^2\\
 &+(18u^5+137u^4+385u^3+484u^2+258u+44)v\\
 &+4u^6+38u^5+140u^4+252u^3+230u^2+104u+24.
\end{align*}
Every coefficient is nonnegative.  This proves
\eqref{det:telescoping}.  Summation in $d$, with
$t_{r,1}=(r-1)/r$, gives
\begin{equation}\label{det:J-bound}
 J_r\leq\frac{2r^3}{r-1}\frac{t_{r,1}}2=r^2.
\end{equation}

For the remaining sum, direct calculation gives
\[
 (I_2,I_3,I_4)=\left(6,4,\frac{25}{4}\right),\qquad
 (J_2,J_3,J_4)=\left(0,6,\frac{43}{4}\right).
\]
Thus \eqref{det:scalar-budget} holds for $r\leq4$.
Also, $I_5=24/5\leq5$ and $I_6=329/60\leq6$.  For $r\geq7$,
$\binom rp\geq\binom r2$ when $2\leq p\leq r/2$, so
\[
 I_r\leq\frac{3(r+1)}r
       +\frac{3(r+1)(r-2)}{r(r-1)}\leq r.
\]
Combine this with \eqref{det:J-bound} to finish the proof.
\end{proof}
